\documentclass[a4paper]{article}
\usepackage[landscape,margin=0.9cm]{geometry}
\usepackage[T1]{fontenc}
\usepackage{lmodern}
\usepackage{tabularx,booktabs,array}
\usepackage{amsmath}
\usepackage{microtype}
\setlength{\parindent}{0pt}
\newcolumntype{Y}{>{\raggedright\arraybackslash}X}
\newcolumntype{P}[1]{>{\raggedright\arraybackslash}p{#1}}
\pagestyle{empty}

\begin{document}
\fontsize{9}{11}\selectfont
\renewcommand{\arraystretch}{1.1}
\setlength{\tabcolsep}{3.5pt}

\noindent
\begin{tabularx}{\textwidth}{@{}P{1.7cm} P{5.6cm} P{6.0cm} P{3.8cm} Y@{}}
\toprule
\textbf{Module} & \textbf{Problem} & \textbf{Candidates} & \textbf{Selection} & \textbf{Justification}\\
\midrule
A.\ Relief inventory &
Input: $n$ crates with weight $w_i$ and value $v_i$ (day-dependent values) and a capacity $W$. Output: a set of indivisible crates of total weight at most $W$ with maximum total value (0/1 knapsack). Tier L: $n\le 1500$, $W\le 21\,000$. &
Brute force $O(2^n)$; greedy by exact ratio $v/w$ $O(n\log n)$ (rule of thumb, not optimal); dynamic programming over the capacity $O(nW)$; branch and bound (worst case $O(2^n)$). &
Primary: dynamic programming $O(nW)$ with backtracking to recover the crates. Baseline: the ratio greedy (ties by crate id), and brute force for $n\le 20$. &
Weights are integers at most 20, so $W\le 0.7\cdot 20n$ and $nW\le 3\cdot 10^{7}$, which makes the exact method affordable. Trap 47 shows that greedy reaches only 55\% of the optimum. The selection fails if weights were very large ($W$ exponential in the input size); branch and bound would then be needed. \emph{Campaign:} loads the truck on days 1--6 with the day's values.\\
\addlinespace
B.\ Route planning &
Input: $V$ locations, $E\le 3V$ roads, directed step cost = minutes + penalty of the destination ($\ge 1$), source HQ, blocked roads. Output: a cheapest route and its cost or ``no route''; the reachable set; a fewest-roads route. &
Dijkstra with binary heap $O((V+E)\log V)$; Dijkstra with array $O(V^2)$; Bellman--Ford $O(VE)$; breadth-first search $O(V+E)$ (number of roads only). &
Primary: heap-based Dijkstra for the cost, and BFS (level by level, neighbours in ascending id) for reachability and fewest roads, on the open roads. Baseline: Bellman--Ford. &
All costs are positive and the graphs are sparse, so heap-based Dijkstra is the best choice. In trap 48 the fewest-roads route is 5.6 times as costly as the cheapest one, so BFS alone is insufficient. The selection fails for negative costs, where Bellman--Ford is required. \emph{Campaign:} checks reachability and chooses the route every day.\\
\addlinespace
C.\ Rescue scheduling &
Input: $n\le 3000$ windows (start, finish, value), some unavailable on a given day; touching windows are compatible. Output: a set of pairwise compatible windows with the maximum number of windows, or with the maximum total value. &
Brute force $O(2^n)$; earliest-finish greedy $O(n\log n)$; weighted-interval dynamic programming with binary search $O(n\log n)$; plain dynamic programming $O(n^2)$. &
Primary: earliest-finish greedy (count) and weighted-interval DP with binary search (value). Baseline: the $O(n^2)$ DP, and brute force on tier S. &
Greedy is optimal for the count (exchange argument) but not for the value, since one valuable long window can block two short ones; the cyclone data overlaps heavily. Compatibility is tested as start $\ge$ previous finish. The method fails if windows also compete for another resource. \emph{Campaign:} staffs the windows that are available that day.\\
\addlinespace
D.\ Emergency queue &
Input: $m\le 5000$ alarms (tick, priority 1--10, id). Output: the release order, one alarm per tick, in urgency order (priority descending, tick ascending, alarm id ascending) among the alarms present. &
Linear scan of the present alarms $O(m\cdot\mathrm{backlog})$; sorted-list insertion $O(m^2)$; binary heap $O(m\log m)$; ten priority buckets $O(m)$. &
Primary: binary heap keyed by (priority descending, tick, alarm id). Baseline: linear scan for the most urgent alarm. &
Rates of 1.2--1.8 alarms per tick against one release per tick make the backlog grow to about $0.4m$, so scanning costs $O(m^2)$ while the heap costs $O(\log m)$ per event. Buckets would also work because there are only ten priorities; the heap is unnecessary only for very small inputs. \emph{Campaign:} processes the alarm stream, including the alarms added on days 2 and 5.\\
\addlinespace
F.\ Network planning &
Input: the roads of the reachable part with cable cost 1--100. Output: roads linking all reachable locations at minimum total cable cost (a spanning tree, or a forest if waves split the map). &
Kruskal with union--find $O(E\log E)$; Prim with heap $O(E\log V)$; Boruvka $O(E\log V)$; Prim with array $O(V^2)$. &
Primary: Kruskal with union--find. Baseline: Prim with a heap. &
Since $E\le 3V$ the graph is sparse, sorting the edges is cheap, and Kruskal returns a forest when the map is disconnected. It is no longer the best choice on dense graphs ($E\approx V^2$), where Prim with an array is preferable.\\
\addlinespace
G.\ Chains and cuts &
Input: dimensions $p_0,\dots,p_n$ ($n\le 300$) and drum prices for lengths $1..L$ ($L\le 2000$). Output: the bracketing with the fewest scalar multiplications ($p\,q\,r$ per product); a cut of the drum of maximum total price. &
Chain: brute force over brackets (Catalan number), memoised recursion $O(n^3)$, bottom-up DP $O(n^3)$. Drum: recursion $O(2^L)$, DP $O(L^2)$; greedy by price per length as rule of thumb. &
Primary: bottom-up DP for both ($n^3\approx 2.7\cdot 10^7$, $L^2\approx 4\cdot 10^6$) with a split table. Baseline: memoised recursion and the greedy cut; brute force on tier S. &
Both problems have optimal substructure and overlapping subproblems, and the prices are not proportional to the length, so the greedy rule is not correct. The cubic method fails for $n\gg 1000$, where the Hu--Shing algorithm would be needed.\\
\bottomrule
\end{tabularx}

\newpage
\fontsize{11}{15}\selectfont
\renewcommand{\arraystretch}{1.5}
\setlength{\tabcolsep}{6pt}
\noindent
\begin{tabularx}{\textwidth}{@{}P{4.4cm} Y@{}}
\toprule
\textbf{Member} & \textbf{Contribution}\\
\midrule
Ayna Khan, 24I-0097, Section C &
\textbf{Iteration 1:} generator foundations, crates module with trap cases 46 and 47, radio module with trap case 49, pipeline, safety rules and rations, the tier-S claims and the workbook writer; validator checks for crates and text and the test runner; approach note Part 3 and assembly; this proposal and the README. \textbf{Iteration 2 proposed:} modules A (relief inventory) and D (emergency queue), and the campaign coordinator.\\
Dur E Adan, 24I-3037, Section C &
\textbf{Iteration 1:} generator network module with trap case 48 (locations, roads, waves, targets), relief sites with the planted tour and the sensor graph; validator route checks and the broken-copy tests including the trap cases; approach note Parts 1 and 4; the members and repository files. \textbf{Iteration 2 proposed:} modules B (route planning) and F (network planning), and the data loader with the validation layer.\\
Hamza Sardar, 24I-0089, Section C &
\textbf{Iteration 1:} generator module for windows and alarms; the validator engine with checks V1--V8 and the variety rules and the exact recomputation of the tier-S claims; approach note Part 2; ran the complete pipeline and produced the cases workbook, the broken-copy workbooks and the validation report. \textbf{Iteration 2 proposed:} modules C (rescue scheduling) and G (chains and cuts), and the experiment harness.\\
\bottomrule
\end{tabularx}

\end{document}
